% ============================================================================
% Agentic Programming: Exercises & Solution Sketches
% ============================================================================

\documentclass[aspectratio=169,10pt]{beamer}

% ---------- Theme and colors ----------
\usecolortheme{dove}
\setbeamertemplate{navigation symbols}{}
\setbeamertemplate{footline}{%
  \hfill\usebeamerfont{page number in head/foot}%
  \usebeamercolor[fg]{normal text}%
  \insertframenumber\,/\,\inserttotalframenumber\kern1em\vskip4pt}
\setbeamerfont{frametitle}{size=\huge}

% ---------- Packages ----------
\usepackage[utf8]{inputenc}
\usepackage[T1]{fontenc}
\usepackage{amsmath,amssymb,amsfonts}
\usepackage{booktabs}
\usepackage{hyperref}
\usepackage{multicol}
\usepackage{xcolor}
\usepackage{listings}
\usepackage{pifont}

% ---------- Listings ----------
\lstset{language=Python,basicstyle=\small\ttfamily,
    keywordstyle=\color{uzhblue}\bfseries,commentstyle=\color{uzhgreydark}\itshape,
    stringstyle=\color{darkgreen},showstringspaces=false,breaklines=true,frame=none,
    xleftmargin=0pt,aboveskip=4pt,belowskip=4pt,}
\lstdefinestyle{bash}{language=bash,basicstyle=\small\ttfamily,
    keywordstyle=\color{uzhblue}\bfseries,commentstyle=\color{uzhgreydark}\itshape,
    showstringspaces=false,breaklines=true,frame=none,
    morekeywords={claude,git,uv,npm},}
\lstdefinestyle{prompt}{basicstyle=\small\ttfamily,breaklines=true,frame=none,
    xleftmargin=6pt,literate={>}{{\textcolor{harvardcrimson}{>}}}1,}

% ---------- Colors ----------
\definecolor{uzhblue}{rgb}{0,0,0.5}
\definecolor{uzhgreylight}{RGB}{236,238,241}
\definecolor{uzhgreydark}{RGB}{81,86,94}
\definecolor{harvardcrimson}{RGB}{153,0,0}
\definecolor{unil@blue}{RGB}{0,140,204}
\definecolor{unil@black}{RGB}{35,31,32}
\definecolor{darkgreen}{RGB}{0,100,0}
\definecolor{darkred}{RGB}{139,0,0}
\definecolor{softblue}{RGB}{70,130,180}
\definecolor{softorange}{RGB}{230,159,0}
\definecolor{softgreen}{RGB}{0,158,115}

\setbeamercolor{frametitle}{bg=uzhgreylight, fg=uzhblue}
\setbeamercolor{title}{fg=uzhblue}
\setbeamercolor{structure}{fg=uzhblue}
\setbeamercolor{block title}{bg=unil@blue, fg=white}
\setbeamercolor{block body}{bg=uzhgreylight}
\setbeamercolor{block title alerted}{bg=darkred,fg=white}
\setbeamercolor{block body alerted}{bg=red!5}
\setbeamercolor{block title example}{bg=darkgreen,fg=white}
\setbeamercolor{block body example}{bg=green!5}
\setbeamercolor{item}{fg=unil@black}
\setbeamercolor{normal text}{fg=unil@black}
\newcommand{\emphc}[1]{\textcolor{harvardcrimson}{#1}}
\hypersetup{colorlinks,linkcolor=harvardcrimson,urlcolor=harvardcrimson}

% ---------- Title ----------
\title[Toolkit T3]{Toolkit T3: Agentic programming, exercise handout}
\subtitle{Workshop and self-study exercises for the agentic-programming toolkit}
\author{Simon Scheidegger}
\institute[HEC Lausanne]{HEC, University of Lausanne}
\date{}

\begin{document}

{\setbeamercolor{background canvas}{bg=uzhgreylight}
\begin{frame}
\titlepage
\begin{center}
{\footnotesize Prompts: \texttt{lectures/lecture\_06\_B17\_agentic\_programming/code/exercise\_prompts.md}\\
Solution notes: \texttt{lectures/lecture\_06\_B17\_agentic\_programming/code/exercise\_solutions.md}\\[2pt]
Materials: \url{https://github.com/sischei/Deep_Learning_for_Solving_And_Estimating_Dynamic_Economic_Models}}
\end{center}
\end{frame}}

% === OVERVIEW ===
\begin{frame}{Exercise Overview}
\scriptsize
\textbf{Classroom track (Ex.\,0--7, 110 min total incl.\ pre-flight):}
\renewcommand{\arraystretch}{1.05}
\begin{center}
\begin{tabular}{c l c l}
\toprule
\textbf{\#} & \textbf{Exercise} & \textbf{Time} & \textbf{Focus} \\
\midrule
0 & Hello Claude (\emphc{pre-flight}) & 10 min & First contact, zero setup \\
1 & Environment Bootstrap & 15 min & Setup \& first Claude session \\
2 & First Real Task & 15 min & Level 3 prompts + Plan Mode \\
3 & Prompt Refactor & 15 min & 6-component prompt framework \\
4 & Write Your Own CLAUDE.md & 10 min & Persistent project memory \\
5 & Skills \& Subagents & 10 min & \texttt{/agents} UI; verifier subagent \\
6 & Data $\to$ Figures $\to$ LaTeX & 25 min & Mini research pipeline \\
7 & Replicate a Method & 10 min & Autonomous paper replication \\
\bottomrule
\end{tabular}
\end{center}

\vspace{0.4em}
\textbf{Self-study extensions (Ex.\,8--12):}
\begin{center}
\begin{tabular}{c l c l}
\toprule
\textbf{\#} & \textbf{Exercise} & \textbf{Time} & \textbf{Focus} \\
\midrule
8  & Mincer Warm-Up Demo & 15--30 min & End-to-end OLS pipeline \\
9  & Specialist Subagent & 20--45 min & Econ / MC / backtest reviewer \\
10 & Hook \& Real Skill & 20--60 min & Hooks + \texttt{/strategic-revision} \\
11 & Monte Carlo Study & 30--45 min & Subagent-driven MC \\
\textbf{12} & \textbf{Ralph Autonomous Loop} & \textbf{30--45 min} & \textbf{Closed-loop autonomy} \\
\bottomrule
\end{tabular}
\end{center}

\vspace{0.2em}
{\tiny Solutions in \texttt{lectures/lecture\_06\_B17\_agentic\_programming/code/exercise\_solutions.md}; consult only \emph{after} attempting.}
\end{frame}

% --- Learning objectives ---
\begin{frame}{Learning objectives}
\textbf{By the end of this lecture you should be able to:}
\begin{itemize}
  \item Compose effective prompts using role, context, task, constraints, output, and verification components
  \item Create reusable skills and subagents to decompose complex analysis workflows
  \item Build end-to-end pipelines from synthetic data generation through estimation, visualization, and LaTeX tables
  \item Use Bayesian active learning and Monte Carlo verification to test agentic reasoning robustness
  \item Replicate published computational results by extracting specifications and delegating multi-step tasks
\end{itemize}
\end{frame}


% ============================================================================
\begin{frame}{}
\vfill\begin{center}{\Huge \textcolor{uzhblue}{\textbf{Exercise 1}}}\\[12pt]
{\LARGE Environment Bootstrap}\\[6pt]{\small 15 minutes}\end{center}\vfill
\end{frame}

\begin{frame}[fragile]{Exercise 1: Environment Bootstrap}
\begin{exampleblock}{Goal}
Get a working Claude Code session in a version-controlled project.
\end{exampleblock}
\small
\begin{enumerate}\itemsep1pt
    \item Create \texttt{\~{}/research/lecture-demo}
    \item Copy in \texttt{generate\_synthetic\_data.py} and run it once
    \item Initialize git and make an initial commit
    \item Launch Claude Code and run \texttt{/init}
    \item Ask Claude to create \texttt{explore\_data.py} with \texttt{pathlib}
    \item Run it on \texttt{data/synthetic\_panel.csv}
\end{enumerate}

\textbf{Success:} local data + working script + \texttt{CLAUDE.md} + git commit
\end{frame}

\begin{frame}{Exercise 1: Solution Sketch}
\small
\begin{exampleblock}{Reference answer}
\begin{itemize}\itemsep2pt
    \item Expected artifacts: \texttt{CLAUDE.md}, \texttt{data/synthetic\_panel.csv}, \texttt{explore\_data.py}
    \item \texttt{explore\_data.py} should accept a CSV path and report shape, columns, dtypes, head, summary stats, and missingness
    \item The file should use \texttt{pathlib.Path}, not \texttt{os.path}
\end{itemize}
\end{exampleblock}

\begin{alertblock}{Common failure}
Students create the script before the data exists locally. Fix that by copying
and running \texttt{generate\_synthetic\_data.py} first.
\end{alertblock}
\end{frame}

% ============================================================================
\begin{frame}{}
\vfill\begin{center}{\Huge \textcolor{uzhblue}{\textbf{Exercise 2}}}\\[12pt]
{\LARGE First Real Task}\\[6pt]{\small 15 minutes}\end{center}\vfill
\end{frame}

\begin{frame}[fragile]{Exercise 2: First Real Task}
\begin{exampleblock}{Goal}
Use the Level 3 prompt pattern on the shared synthetic panel.
\end{exampleblock}
\small
\begin{lstlisting}[style=prompt, basicstyle=\scriptsize\ttfamily]
> I have a dataset at data/synthetic_panel.csv.
  Without loading it yet, first tell me how to inspect it.
  Then:
  1. Load it and print diagnostics
  2. Compute annual average revenue by treatment group
  3. Generate a 2x2 panel of plots
  4. Save the figure to figures/panel_overview.pdf
  5. Write a 3-sentence description to notes/data_description.md
\end{lstlisting}

\textbf{Reflect:} Where did Claude excel? Where did it need correction?
\end{frame}

\begin{frame}{Exercise 2: Solution Sketch}
\small
\begin{columns}[T]
\column{0.48\textwidth}
\begin{exampleblock}{Expected artifacts}
\begin{itemize}\itemsep1pt
    \item \texttt{figures/panel\_overview.pdf}
    \item \texttt{notes/data\_description.md}
    \item A short diagnostic printout in the terminal
\end{itemize}
\end{exampleblock}

\column{0.48\textwidth}
\begin{alertblock}{Typical interpretation}
\begin{itemize}\itemsep1pt
    \item Revenue trends are similar pre-2015
    \item Post-2015 treated firms drift modestly upward
    \item Log revenue and log employees are positively related
\end{itemize}
\end{alertblock}
\end{columns}
\end{frame}

% ============================================================================
\begin{frame}{}
\vfill\begin{center}{\Huge \textcolor{uzhblue}{\textbf{Exercise 3}}}\\[12pt]
{\LARGE Prompt Refactor}\\[6pt]{\small 15 minutes}\end{center}\vfill
\end{frame}

\begin{frame}[fragile]{Exercise 3: Prompt Refactor}
\begin{exampleblock}{Goal}
Transform a bad prompt using the 6-component framework.
\end{exampleblock}
\small
\textbf{Bad prompt:} \texttt{> Analyze the data and run the regression and make a table.}

\vspace{0.3em}
\textbf{Task:}
\begin{enumerate}\itemsep1pt
    \item Diagnose the missing components
    \item Rewrite as Level 3
    \item Compare the vague and rewritten versions in Claude Code
\end{enumerate}
\end{frame}

\begin{frame}[fragile]{Exercise 3: Solution Sketch}
\scriptsize
\begin{lstlisting}[style=prompt, basicstyle=\tiny\ttfamily]
> [ROLE] You are a computational economist working on panel data analysis.
  [CONTEXT] I have a firm-level panel dataset at data/synthetic_panel.csv.
  Variables: revenue, employees, treated, firm_id, year. Treatment starts in 2015.
  [TASK] Estimate log(revenue) ~ treated * post + firm_FE + year_FE
  with standard errors clustered at the firm level.
  [CONSTRAINTS] Use pyfixest. Do not modify the data file.
  [OUTPUT] Save regression results to results/did_estimates.json
  and a LaTeX table to paper/tables/tab_did.tex.
  [VERIFY] Also test on synthetic data where true effect = 0.05.
\end{lstlisting}

\begin{block}{Key}
Every component is explicit. Two different students should be able to read this
prompt and ask Claude for essentially the same deliverable.
\end{block}
\end{frame}

% ============================================================================
\begin{frame}{}
\vfill\begin{center}{\Huge \textcolor{uzhblue}{\textbf{Exercise 4}}}\\[12pt]
{\LARGE Write Your Own CLAUDE.md}\\[6pt]{\small 10 minutes}\end{center}\vfill
\end{frame}

\begin{frame}[fragile]{Exercise 4: Write Your Own CLAUDE.md}
\begin{exampleblock}{Goal}
Create a project memory file and test both normal behavior and guardrails.
\end{exampleblock}
\small
\begin{enumerate}\itemsep2pt
    \item Copy \texttt{lectures/lecture\_06\_B17\_agentic\_programming/code/CLAUDE\_md\_template.md}
    \item Fill in Repository Structure, Coding Standards, and Current Status
    \item Start a fresh Claude session and ask for a project summary
    \item Test one normal coding request
    \item Test one forbidden request touching a read-only path
\end{enumerate}
\end{frame}

\begin{frame}{Exercise 4: Solution Sketch}
\small
\begin{exampleblock}{What success looks like}
\begin{itemize}\itemsep2pt
    \item Claude uses your notation, folder structure, and coding rules immediately
    \item Claude respects \texttt{pathlib}, type hints, and naming conventions
    \item Claude refuses or pushes back on edits to \texttt{data/raw/}
\end{itemize}
\end{exampleblock}

\begin{alertblock}{High-value sections}
\texttt{Repository Structure}, \texttt{Coding Standards}, and
\texttt{Current Status} have the highest immediate payoff.
\end{alertblock}
\end{frame}

% ============================================================================
\begin{frame}{}
\vfill\begin{center}{\Huge \textcolor{uzhblue}{\textbf{Exercise 5}}}\\[12pt]
{\LARGE Skills \& Subagents}\\[6pt]{\small 10 minutes}\end{center}\vfill
\end{frame}

\begin{frame}[fragile]{Exercise 5: Create a Skill \& Use Subagents}
\begin{exampleblock}{Goal}
Create one reusable skill and one verifier-style subagent.
\end{exampleblock}
\small
\begin{columns}[T]
\column{0.48\textwidth}
\textbf{Part A: Skill}
\begin{enumerate}\itemsep1pt
    \item Copy \texttt{lectures/lecture\_06\_B17\_agentic\_programming/code/skills/example\_skill/SKILL.md}
    \item Save it to \texttt{.claude/skills/data-diagnostics/}\\
    \texttt{SKILL.md}
    \item Run \texttt{/data-diagnostics data/synthetic\_panel.csv}
\end{enumerate}

\column{0.48\textwidth}
\textbf{Part B: Subagent}
\begin{enumerate}\itemsep1pt
    \item Copy \texttt{lectures/lecture\_06\_B17\_agentic\_programming/code/subagents/}\\
    \texttt{verifier.md}
    \item Save it to \texttt{.claude/agents/verifier.md}
    \item Ask Claude to use the verifier subagent on the diagnostics report
\end{enumerate}
\end{columns}
\end{frame}

\begin{frame}{Exercise 5: Solution Sketch}
\small
\begin{columns}[T]
\column{0.48\textwidth}
\begin{exampleblock}{Expected artifacts}
\begin{itemize}\itemsep1pt
    \item \texttt{.claude/skills/data-diagnostics/}\\
    \texttt{SKILL.md}
    \item \texttt{.claude/agents/verifier.md}
    \item Diagnostics note in \texttt{notes/}
\end{itemize}
\end{exampleblock}

\column{0.48\textwidth}
\begin{alertblock}{Conceptual distinction}
\begin{itemize}\itemsep1pt
    \item \textbf{Skill:} repeatable multi-step workflow
    \item \textbf{Subagent:} isolated verifier or researcher
    \item Use subagents when you want context separation
\end{itemize}
\end{alertblock}
\end{columns}
\end{frame}

% ============================================================================
\begin{frame}{}
\vfill\begin{center}{\Huge \textcolor{uzhblue}{\textbf{Exercise 6}}}\\[12pt]
{\LARGE Data $\to$ Figures $\to$ LaTeX}\\[6pt]{\small 25 minutes}\end{center}\vfill
\end{frame}

\begin{frame}[fragile]{Exercise 6: Build a Mini Pipeline}
\begin{exampleblock}{Goal}
Build a pipeline from shared synthetic data to a publication-ready figure and
LaTeX table.
\end{exampleblock}
\scriptsize
\begin{enumerate}\itemsep2pt
    \item Run \texttt{/data-diagnostics data/synthetic\_panel.csv}
    \item Write \texttt{code/estimate\_did.py}
    \item Write \texttt{code/plot\_did.py}
    \item Write \texttt{code/make\_table.py}
    \item Verify the estimated effect against the true value 0.05
\end{enumerate}

\textbf{Stretch:} add \texttt{run\_all.py} with a \texttt{--dry-run} flag.
\end{frame}

\begin{frame}{Exercise 6: Solution Sketch}
\small
\begin{exampleblock}{Expected artifacts}
\begin{itemize}\itemsep1pt
    \item \texttt{results/did\_results.json}
    \item \texttt{figures/event\_study.pdf} and \texttt{.png}
    \item \texttt{paper/tables/tab\_did.tex}
    \item Optional \texttt{run\_all.py}
\end{itemize}
\end{exampleblock}

\begin{alertblock}{Minimum verification}
\begin{itemize}\itemsep1pt
    \item Two-way FE coefficient near 0.05
    \item Statistically significant at the 5\% level
    \item Pre-trend coefficients near zero
\end{itemize}
\end{alertblock}
\end{frame}

% ============================================================================
\begin{frame}{}
\vfill\begin{center}{\Huge \textcolor{uzhblue}{\textbf{Exercise 7}}}\\[12pt]
{\LARGE Replicate a Method}\\[6pt]{\small 10 minutes}\end{center}\vfill
\end{frame}

\begin{frame}[fragile]{Exercise 7: Paper Replication}
\begin{exampleblock}{Goal}
Use an AI agent to replicate a published result.
\end{exampleblock}
\small
\textbf{Choose:} \textcolor{darkgreen}{\textbf{Beginner}} (Lasso, Stock\&Watson 2012) $|$
\textcolor{softorange}{\textbf{Intermediate}} (SINDy, Brunton 2016) $|$
\textcolor{harvardcrimson}{\textbf{Advanced}} (DEQN from earlier lectures)

\vspace{0.3em}
\begin{lstlisting}[style=prompt, basicstyle=\scriptsize\ttfamily]
> Role: Replication engineer.
  Paper: [title]. Mission: Replicate [figure/table].
  Rules: no author code; track notes/todo.md; commit after each step.
  Before implementing: extract figure axes, ranges, and CIs.
  Write plan to notes/plan.md. Ask me to confirm.
\end{lstlisting}
\end{frame}

\begin{frame}{Exercise 7: Solution Sketch}
\small
\begin{exampleblock}{Success criteria}
\begin{itemize}\itemsep2pt
    \item A bounded target figure or table
    \item Progress tracked in \texttt{notes/plan.md} and \texttt{notes/todo.md}
    \item Final comparison written to \texttt{notes/replication\_assessment.md}
\end{itemize}
\end{exampleblock}

\begin{alertblock}{Stop rule}
If the agent is stuck after three serious attempts, write
\texttt{notes/blocked.md} and stop. ``The code runs'' is not the same as
``the result matches.''
\end{alertblock}
\end{frame}

% ============================================================================
\begin{frame}{}
\vfill\begin{center}{\Huge \textcolor{uzhblue}{\textbf{Appendix}}}\\[12pt]
{\LARGE Essential Prompt Library}\end{center}\vfill
\end{frame}

\begin{frame}[fragile]{Prompt Library}
\scriptsize
\begin{columns}[T]
\column{0.48\textwidth}
\textbf{Session:}
\begin{lstlisting}[style=prompt, basicstyle=\tiny\ttfamily]
> Write notes/session_YYYY-MM-DD.md:
  what we did, worked, failed, next.
\end{lstlisting}
\textbf{Code Review:}
\begin{lstlisting}[style=prompt, basicstyle=\tiny\ttfamily]
> Review .py files. Refactor >50-line
  functions. Add missing docstrings.
\end{lstlisting}
\textbf{Tests:}
\begin{lstlisting}[style=prompt, basicstyle=\tiny\ttfamily]
> Write pytest tests with known-answer
  synthetic data.
\end{lstlisting}
\textbf{Debug:}
\begin{lstlisting}[style=prompt, basicstyle=\tiny\ttfamily]
> Error: [paste]. Don't fix yet.
  Explain cause. Propose 2 fixes.
\end{lstlisting}

\column{0.48\textwidth}
\textbf{Data:}
\begin{lstlisting}[style=prompt, basicstyle=\tiny\ttfamily]
> Load [file]. Do NOT modify.
  Print shape, dtypes, missing%,
  summary stats. Flag outliers.
\end{lstlisting}
\textbf{Merge:}
\begin{lstlisting}[style=prompt, basicstyle=\tiny\ttfamily]
> Left join [A] with [B] on [key].
  Assert no rows dropped.
\end{lstlisting}
\textbf{Writing:}
\begin{lstlisting}[style=prompt, basicstyle=\tiny\ttfamily]
> Improve clarity. Don't change
  meaning or citations.
\end{lstlisting}
\textbf{Citations:}
\begin{lstlisting}[style=prompt, basicstyle=\tiny\ttfamily]
> Do NOT invent citations.
  Only cite papers in references.bib.
\end{lstlisting}
\end{columns}
\end{frame}

\end{document}
